\begingroup
\section{Foundation F65: the exact angular disk endpoint}
\label{module:F65}
\noindent\textit{Audited source: \nolinkurl{convex_largeN_polya_szego_stage65_correct_source_disk_endpoint.tex}.}
\subsection{Fan, correct source, and the two support spaces}

Let
\begin{equation}\label{F65:eq:fan}
 \alpha_i>0,\qquad \sum_i\alpha_i=2\pi,
 \qquad \theta_{i+1}-\theta_i=\alpha_i,
 \qquad \bar\alpha=\frac{2\pi}{N},
\end{equation}
and put
\begin{equation}\label{F65:eq:SJ}
 S=\alpha_{\max}+\bar\alpha,
 \qquad
 \Jfour(\alpha)=\sum_i\alpha_i^4-N\bar\alpha^4.
\end{equation}
Set $\phi_i=\theta_i+\alpha_i/2$ and
$K_i=[\phi_{i-1},\phi_i]$.  On $K_i$ write
\begin{equation}
 a=\alpha_{i-1},\quad b=\alpha_i,\quad
 x=\theta-\theta_i,\quad
 s=\frac{x+a/2}{(a+b)/2},
\end{equation}
and define the correct principal source
\begin{equation}\label{F65:eq:G}
 G_\alpha(\theta)=\frac{x^2}{2}.
\end{equation}

For a coefficient vector $z$, put
\begin{equation}
 r_- =\frac{z_i-z_{i-1}}a,
 \qquad r_+=\frac{z_{i+1}-z_i}b.
\end{equation}
The homogeneous and exact vertex-compatible traces are
\begin{align}
 (T_\alpha^0z)_i(x)
 &=z_i+x\bigl((1-s)r_-+sr_+\bigr),\label{F65:eq:T0}\\
 (T_\alpha z)_i(x)
 &=z_i\cos x+\bigl((1-s)a_i^-+sa_i^+\bigr)\sin x,
 \label{F65:eq:T}
\end{align}
where
\begin{equation}\label{F65:eq:apm}
 a_i^-=r_-\frac a{\sin a}-z_i\tan\frac a2,
 \qquad
 a_i^+=r_+\frac b{\sin b}+z_i\tan\frac b2.
\end{equation}
Both traces are continuous at the vertices $\phi_i$.

The homogeneous and exact area rows are
\begin{align}
 L_\alpha^0z
 &=\frac12\sum_i(\alpha_{i-1}+\alpha_i)z_i,\label{F65:eq:L0}\\
 L_\alpha z
 &=\sum_i\left(\tan\frac{\alpha_{i-1}}2
 +\tan\frac{\alpha_i}2\right)z_i.\label{F65:eq:L}
\end{align}
Let $Z_\alpha^0=\ker L_\alpha^0$ and $Z_\alpha=\ker L_\alpha$.
Centering rows may be added throughout; translations will be removed exactly
in Section~\ref{F65:sec:translations}.

Work in $H=H^{1/2}(\Sone)/\R$ with
\begin{equation}\label{F65:eq:qplus}
 \qplus(f)=\sum_{n\ne0}(|n|+1)|\widehat f_n|^2,
 \qquad
 \widehat f_n=\frac1{2\pi}\int_0^{2\pi}f(\theta)e^{-in\theta}\,d\theta.
\end{equation}
Set
\begin{equation}
 V_\alpha^0=T_\alpha^0Z_\alpha^0,
 \qquad V_\alpha=T_\alpha Z_\alpha,
\end{equation}
and
\begin{align}
 P_0^G(\alpha)&=\inf_{v\in V_\alpha^0}\qplus(G_\alpha+v),
 \label{F65:eq:P0}\\
 P_{\rm tr}^G(\alpha)&=\inf_{v\in V_\alpha}\qplus(G_\alpha+v).
 \label{F65:eq:Ptr}
\end{align}
The corrected homogeneous theorem gives
\begin{equation}\label{F65:eq:input-gap}
 P_0^G(\alpha)-P_0^G(\bar\alpha\mathbf1)
 \ge\frac{\zeta(3)}{10\pi^4}\Jfour(\alpha).
\end{equation}

\subsection{No-ratio exact-trace graph}

\begin{lemma}[Local jet coercivity]\label{F65:lem:jet}
There are absolute constants $0<c<C$ such that, with $w=a+b$,
\begin{equation}\label{F65:eq:jet}
 cw\bigl(z_i^2+w^2r_-^2+w^2r_+^2\bigr)
 \le\|T_\alpha^0z\|_{L^2(K_i)}^2
 \le Cw\bigl(z_i^2+w^2r_-^2+w^2r_+^2\bigr).
\end{equation}
The constants are independent of $a/b$.
\end{lemma}

\begin{proof}
After $x=(w/2)(y-t)$ with $t=a/w$, the three normalized polynomials are
\begin{equation}
 1,\qquad (y-t)(1-y),\qquad (y-t)y,
 \qquad 0\le y\le1.
\end{equation}
The matrix carrying the two slope variables to the linear and quadratic
coefficients has determinant one and it and its inverse are uniformly
bounded for $0\le t\le1$.  Finite-dimensional norm equivalence proves
\eqref{F65:eq:jet}.
\end{proof}

\begin{lemma}[Exact-minus-homogeneous graph]\label{F65:lem:graph}
For $S$ sufficiently small there is a bijection
$C_\alpha:Z_\alpha^0\to Z_\alpha$ for which
\begin{equation}
 J_\alpha(T_\alpha^0z)=T_\alpha C_\alpha z
\end{equation}
satisfies
\begin{equation}\label{F65:eq:Jclose}
 \|(J_\alpha-I)v\|_H\le CS^{3/2}\|v\|_H,
 \qquad v\in V_\alpha^0.
\end{equation}
\end{lemma}

\begin{proof}
Taylor's formula in \eqref{F65:eq:T}--\eqref{F65:eq:apm} gives on $K_i$
\begin{align}
 \|T_\alpha z-T_\alpha^0z\|_{L^\infty(K_i)}
 &\le C\bigl(w^2|z_i|+w^3(|r_-|+|r_+|)\bigr),\notag\\
 \|\partial_\theta(T_\alpha z-T_\alpha^0z)\|_{L^\infty(K_i)}
 &\le C\bigl(w|z_i|+w^2(|r_-|+|r_+|)\bigr).
 \label{F65:eq:trace-jets}
\end{align}
The error is continuous at every wall.  Summing \eqref{F65:eq:trace-jets},
using Lemma~\ref{F65:lem:jet}, cyclic Poincar\'e on the area kernel, and
$\|f\|_{H^{1/2}}^2\le C\|f\|_{L^2}\|f\|_{H^1}$ gives
\begin{equation}\label{F65:eq:trace-close}
 \|T_\alpha z-T_\alpha^0z\|_H
 \le CS^{3/2}\|T_\alpha^0z\|_H.
\end{equation}

For $z\in Z_\alpha^0$, define
\begin{equation}
 C_\alpha z=z-\frac{L_\alpha z}{L_\alpha\mathbf1}\mathbf1.
\end{equation}
Since
\begin{equation}
 \left|\tan\frac u2-\frac u2\right|\le Cu^3,
\end{equation}
the coefficient correction is $O(S^2)$ in the trace metric.  The analogous
formula with $L_\alpha^0$ is the inverse map.  Combining this row correction
with \eqref{F65:eq:trace-close} proves \eqref{F65:eq:Jclose}.
\end{proof}

\begin{lemma}[Hilbert Schur perturbation]\label{F65:lem:schur}
If $J:V_0\to V_1$ is a bijection with
$\|(J-I)v\|\le\eta\|v\|$, $\eta<1/2$, then
\begin{equation}\label{F65:eq:schur}
 \left|\operatorname {dist}(b,V_1)^2
 -\operatorname {dist}(b,V_0)^2\right|
 \le C\eta\|b\|^2.
\end{equation}
\end{lemma}

\begin{proof}
Map the orthogonal projection of $-b$ onto $V_0$ through $J$ and use it as
a competitor in $V_1$.  Apply the same argument to $J^{-1}$.
\end{proof}

The continuous edge quadratic satisfies the no-ratio spline bound
\begin{equation}\label{F65:eq:Gbound}
 \qplus(G_\alpha)\le C\sum_i\alpha_i^4\le CS^3.
\end{equation}
Indeed, affine normalization on each $K_i$ controls the same-cell and
adjacent-cell terms, while dyadic annuli control the separated-cell part.
Lemma~\ref{F65:lem:schur} therefore yields, for
\begin{equation}
 R_{\rm tr}^G=P_{\rm tr}^G-P_0^G,
\end{equation}
the absolute estimate
\begin{equation}\label{F65:eq:trace-absolute}
 |R_{\rm tr}^G(\alpha)|\le CS^{9/2}.
\end{equation}

\subsection{Relative trace comparison}

\begin{proposition}[Quasiuniform trace Hessian]\label{F65:prop:trace-hessian}
Suppose
\begin{equation}\label{F65:eq:quasiuniform}
 (1-\eta_0)a\le\gamma_i\le(1+\eta_0)a,
 \qquad \sum_i\gamma_i=2\pi,
\end{equation}
where $0<\eta_0<1/4$.  If $\sum_i\delta_i=0$, then
\begin{equation}\label{F65:eq:trace-hessian}
 |D^2R_{\rm tr}^G(\gamma)[\delta,\delta]|
 \le Ca^{1/2}\sum_i\gamma_i^2\delta_i^2.
\end{equation}
\end{proposition}

\begin{proof}
If $\delta\ne0$, complexify
$\gamma(z)=\gamma+z\delta$ in the disk
\begin{equation}
 |z|<r,\qquad r=\frac{c_0a}{\|\delta\|_{\ell^\infty}}.
\end{equation}
For every cyclic arc $I$,
\begin{equation}
 \left|z\sum_{i\in I}\delta_i\right|
 \le c_0a\#I
 \le Cc_0\sum_{i\in I}\gamma_i.
\end{equation}
Thus the cellwise affine material map changes every arc length by a fixed
small relative amount.  The pulled-back Gagliardo kernel, the $L^2$ row,
the two traces, the area graph, and the polynomial source $G_{\gamma(z)}$
are holomorphic.  Their Gram inverses stay uniformly bounded by
Lemma~\ref{F65:lem:jet}.  The proof of \eqref{F65:eq:trace-absolute} applies to the
complex bilinear Schur formula and gives
\begin{equation}
 |R_{\rm tr}^G(\gamma(z))|\le Ca^{9/2}.
\end{equation}
Cauchy's estimate gives
\begin{equation}
 |D^2R_{\rm tr}^G(\gamma)[\delta,\delta]|
 \le Ca^{5/2}\|\delta\|_{\ell^\infty}^2.
\end{equation}
Use
$\sum_i\gamma_i^2\delta_i^2\ge c a^2
\|\delta\|_{\ell^\infty}^2$ to obtain \eqref{F65:eq:trace-hessian}.
\end{proof}

\begin{lemma}[Relative trace remainder]\label{F65:lem:trace-relative}
For every sufficiently fine positive fan,
\begin{equation}\label{F65:eq:trace-relative}
 |R_{\rm tr}^G(\alpha)-R_{\rm tr}^G(\bar\alpha\mathbf1)|
 \le C(S^{1/4}+S^{1/2})\Jfour(\alpha).
\end{equation}
\end{lemma}

\begin{proof}
Put $a=\bar\alpha$ and $d_i=\alpha_i-a$.  The exact identity
\begin{equation}\label{F65:eq:Jidentity}
 \Jfour(\alpha)
 =\sum_i d_i^2(\alpha_i^2+2a\alpha_i+3a^2)
\end{equation}
implies $\max_i|d_i|^4\le C\Jfour(\alpha)$.
If $\Jfour\ge S^{17/4}$, use \eqref{F65:eq:trace-absolute}.  Otherwise
$\max|d_i|/S\le CS^{1/16}$, so the segment
$a\mathbf1+td$ is quasiuniform.  Cyclic invariance makes the first
derivative of $R_{\rm tr}^G$ vanish at $a\mathbf1$ on the zero-sum
direction.  Taylor's formula, Proposition~\ref{F65:prop:trace-hessian}, and
\eqref{F65:eq:Jidentity} give the $CS^{1/2}\Jfour$ bound.  The two cases prove
\eqref{F65:eq:trace-relative}.
\end{proof}

\begin{theorem}[Correct-source exact-trace gap]\label{F65:thm:trace-gap}
There is $S_{\rm tr}>0$ such that, for $S\le S_{\rm tr}$,
\begin{equation}\label{F65:eq:trace-gap}
 \boxed{
 P_{\rm tr}^G(\alpha)-P_{\rm tr}^G(\bar\alpha\mathbf1)
 \ge\frac{\zeta(3)}{20\pi^4}\Jfour(\alpha).}
\end{equation}
\end{theorem}

\begin{proof}
Combine \eqref{F65:eq:input-gap} and \eqref{F65:eq:trace-relative}, decreasing
$S_{\rm tr}$ until the latter coefficient is at most
$\zeta(3)/(20\pi^4)$.
\end{proof}

\subsection{Exact translations}\label{F65:sec:translations}

\begin{lemma}[Translation quotient]\label{F65:lem:translations}
For $c\in\R^2$, let $z_i=c\cdot(\cos\theta_i,\sin\theta_i)$.  Then
\begin{equation}\label{F65:eq:translation}
 L_\alpha z=0,
 \qquad T_\alpha z(\theta)=c\cdot(\cos\theta,\sin\theta).
\end{equation}
Consequently
\begin{equation}\label{F65:eq:translation-quotient}
 P_{\rm tr}^G(\alpha)
 =\inf_{v\in\PiNG V_\alpha}
 \qplus(\PiNG G_\alpha+v),
\end{equation}
where $\PiNG$ removes the modes $0,\pm1$.
\end{lemma}

\begin{proof}
With $\tau_i=(-\sin\theta_i,\cos\theta_i)$, the identities
\begin{equation}
 \nu_{i-1}=\nu_i\cos a-\tau_i\sin a,
 \qquad
 \nu_{i+1}=\nu_i\cos b+\tau_i\sin b
\end{equation}
give $a_i^-=a_i^+=c\cdot\tau_i$ in \eqref{F65:eq:apm}.  Hence
\eqref{F65:eq:T} equals $c\cdot\nu(\theta)$.  Translation invariance of area
gives $L_\alpha z=0$.  Orthogonality of the Fourier modes proves
\eqref{F65:eq:translation-quotient}.
\end{proof}

\subsection{Bessel order-minus-one reduction}

Let $j=j_{0,1}$ and, for $n\ge2$, set
\begin{equation}
 k_n=j\frac{J_{n+1}(j)}{J_n(j)}.
\end{equation}
On the range of $\PiNG$, define
\begin{align}
 k(f,g)&=\sum_{|n|\ge2}k_{|n|}\widehat f_n\widehat g_{-n},\notag\\
 \qdisk(f,g)&=\qplus(f,g)-k(f,g).
 \label{F65:eq:qdisk}
\end{align}
The Bessel recurrence and its alternating fixed-argument series give
\begin{equation}\label{F65:eq:bessel-bounds}
 0<k_n<\frac4n,
 \qquad \frac13\qplus(f)\le\qdisk(f)\le\qplus(f).
\end{equation}
Put
\begin{equation}
 b_\alpha^G=\PiNG G_\alpha,
 \qquad W_\alpha=\PiNG V_\alpha,
\end{equation}
and
\begin{equation}\label{F65:eq:PDG}
 P_D^G(\alpha)=\inf_{v\in W_\alpha}\qdisk(b_\alpha^G+v).
\end{equation}

\begin{lemma}[No-ratio exact-trace quasi-interpolant]
\label{F65:lem:QI}
There is a linear operator
\begin{equation}
 I_\alpha:H_{\rm ng}^{3/2}(\Sone)\longrightarrow W_\alpha
\end{equation}
such that, for every sufficiently fine positive fan,
\begin{equation}\label{F65:eq:interpolant}
 \|\psi-I_\alpha\psi\|_{H^{1/2}}
 \le CS\|\psi\|_{H^{3/2}}.
\end{equation}
The constant is independent of the smallest cell and every adjacent
mesh ratio.
\end{lemma}

\begin{proof}
It is enough to work first in the homogeneous trace space.  Put
\(M=\lfloor S^{-1}\rfloor\) and let \(\phi=P_{\le M}\psi\) be the
Fourier projection onto \(|n|\le M\).  Directly from the Fourier norms,
\begin{equation}\label{F65:eq:QI-tail}
 \|\psi-\phi\|_{H^{1/2}}
 \le CS\|\psi\|_{H^{3/2}},
 \qquad
 \|\phi''\|_{L^2}
 \le CS^{-1/2}\|\psi\|_{H^{3/2}}.
\end{equation}

Take the nodal values \(z_i=\phi(\theta_i)\) and set
\(Q_\alpha\phi=T_\alpha^0z\).  We give the global estimate rather than
invoke a mesh-regular interpolation theorem.  On
\(K_i=[\phi_{i-1},\phi_i]\), write
\(a=\alpha_{i-1}\), \(b=\alpha_i\), \(w=a+b\), and
\(\omega_i=[\theta_{i-1},\theta_{i+1}]\).  Since
\begin{equation*}
 r_-=\frac1a\int_{\theta_{i-1}}^{\theta_i}\phi'(u)\,du,
 \qquad
 r_+=\frac1b\int_{\theta_i}^{\theta_{i+1}}\phi'(u)\,du,
\end{equation*}
Cauchy--Schwarz and the fundamental theorem of calculus give
\begin{equation}\label{F65:eq:QI-slope}
 |r_--\phi'(\theta_i)|+|r_+-\phi'(\theta_i)|
 \le Cw^{1/2}\|\phi''\|_{L^2(\omega_i)}.
\end{equation}
The trace formula \eqref{F65:eq:T0} reproduces every affine function.  Using
\eqref{F65:eq:QI-slope} in that formula, and using Taylor's integral formula
for \(\phi\), yields for \(x\in K_i\)
\begin{align}
 |\phi(x)-Q_\alpha\phi(x)|
 &\le Cw^{3/2}\|\phi''\|_{L^2(\omega_i)},\label{F65:eq:QI-point}\\
 |\phi'(x)-(Q_\alpha\phi)'(x)|
 &\le Cw^{1/2}\|\phi''\|_{L^2(\omega_i)}.\label{F65:eq:QI-derivative}
\end{align}
Indeed, after subtracting the affine Taylor polynomial at \(\theta_i\),
the first row consists of \(x\) times a convex combination of the two
slope errors, and the differentiated row contains in addition
\(2x(r_+-r_-)/w\); all its coefficients are bounded for
\(-a/2\le x\le b/2\), including when \(a/b\) tends to zero or infinity.

The intervals \(\omega_i\) have overlap at most two when their
\(L^2\)-mass is assigned back to the two constituent normal cells.
Since \(w\le2S\), integration of
\eqref{F65:eq:QI-point}--\eqref{F65:eq:QI-derivative} and summation give
\begin{equation}\label{F65:eq:QI-L2H1}
 \|\phi-Q_\alpha\phi\|_{L^2}
 \le CS^2\|\phi''\|_{L^2},
 \qquad
 \|\phi-Q_\alpha\phi\|_{H^1}
 \le CS\|\phi''\|_{L^2}.
\end{equation}
The Fourier interpolation inequality
\(\|e\|_{H^{1/2}}^2\le\|e\|_{L^2}\|e\|_{H^1}\),
\eqref{F65:eq:QI-tail}, and \eqref{F65:eq:QI-L2H1} imply
\begin{equation}\label{F65:eq:QI-homogeneous}
 \|\phi-Q_\alpha\phi\|_{H^{1/2}}
 \le CS^{3/2}\|\phi''\|_{L^2}
 \le CS\|\psi\|_{H^{3/2}}.
\end{equation}

It remains only to impose the exact finite rows.  Replace \(z\) by
\begin{equation*}
 z^0=z-\frac{L_\alpha^0z}{L_\alpha^0\mathbf1}\mathbf1.
\end{equation*}
Because \(T_\alpha^0\mathbf1=1\), this changes the homogeneous trace only
by a constant.  Put \(v^0=T_\alpha^0z^0\in V_\alpha^0\) and define
\begin{equation}
 I_\alpha\psi=\PiNG J_\alpha v^0\in W_\alpha,
\end{equation}
where \(J_\alpha\) is the exact area graph of
Lemma~\ref{F65:lem:graph}.  The projection \(\PiNG\) is an orthogonal
contraction in every Fourier Sobolev norm, \(\psi=\PiNG\psi\), and
\eqref{F65:eq:Jclose} gives
\begin{equation*}
 \|\PiNG(J_\alpha-I)v^0\|_{H^{1/2}}
 \le CS^{3/2}\|v^0\|_{H^{1/2}}
 \le CS\|\psi\|_{H^{3/2}}.
\end{equation*}
Together with \eqref{F65:eq:QI-tail} and
\eqref{F65:eq:QI-homogeneous}, this is \eqref{F65:eq:interpolant}.
\end{proof}

\begin{lemma}[Galerkin negative-norm gain]\label{F65:lem:negative}
Let $f_\alpha=b_\alpha^G+v_\alpha$ be the $\qplus$-minimizing residual.
Then
\begin{equation}\label{F65:eq:negative}
 \|f_\alpha\|_{H^{-1/2}}
 \le CS\|f_\alpha\|_{H^{1/2}}.
\end{equation}
\end{lemma}

\begin{proof}
For $g\in H_{\rm ng}^{1/2}$ solve $(|D|+I)\psi=g$.  Galerkin
orthogonality and Lemma~\ref{F65:lem:QI} give
\begin{equation}
 |\langle f_\alpha,g\rangle|
 =|\qplus(f_\alpha,\psi-I_\alpha\psi)|
 \le CS\|f_\alpha\|_{H^{1/2}}\|g\|_{H^{1/2}}.
\end{equation}
Taking the dual supremum proves \eqref{F65:eq:negative}.
\end{proof}

Let
\begin{equation}
 R_B^G(\alpha)=P_{\rm tr}^G(\alpha)-P_D^G(\alpha).
\end{equation}
Completing the disk square gives the exact identity
\begin{equation}\label{F65:eq:bessel-schur}
 R_B^G
 =k(f_\alpha,f_\alpha)
 +\langle g_\alpha,A_{D,\alpha}^{-1}g_\alpha\rangle,
 \qquad g_\alpha(v)=k(f_\alpha,v),
\end{equation}
where $A_{D,\alpha}=\qdisk|_{W_\alpha}$.
Equations \eqref{F65:eq:bessel-bounds}, \eqref{F65:eq:negative},
\eqref{F65:eq:Gbound}, and \eqref{F65:eq:bessel-schur} imply
\begin{equation}\label{F65:eq:bessel-absolute}
 0\le R_B^G(\alpha)\le CS^5.
\end{equation}

\begin{proposition}[Relative Bessel correction]\label{F65:prop:bessel-relative}
For every sufficiently fine positive fan,
\begin{equation}\label{F65:eq:bessel-relative}
 |R_B^G(\alpha)-R_B^G(\bar\alpha\mathbf1)|
 \le C(S^{1/2}+S)\Jfour(\alpha).
\end{equation}
\end{proposition}

\begin{proof}
In a quasiuniform band, pull all functions to a fixed real fan.  If
$h=\|\delta\|_{\ell^\infty}/a$, the first two material derivatives of
$\qplus$ are bounded by $Ch^r$ in $H^{1/2}$, while those of $k$ are bounded
by $Ch^r$ in $H^{-1/2}$.  The latter follows from
\begin{equation}
 |\Delta^rk_n|\le C_r(1+n)^{-1-r},\qquad 0\le r\le3,
\end{equation}
and three discrete summations by parts in the pulled-back kernel.
Differentiating the Galerkin equation and \eqref{F65:eq:negative} twice gives
\begin{equation}
 \|\partial_t^rf_t\|_{H^{1/2}}\le Ch^ra^{3/2},
 \qquad
 \|\partial_t^rf_t\|_{H^{-1/2}}\le Ch^ra^{5/2},
 \quad r=0,1,2.
\end{equation}
Twice differentiating \eqref{F65:eq:bessel-schur} therefore yields
\begin{equation}\label{F65:eq:bessel-hessian}
 |D^2R_B^G(\gamma)[\delta,\delta]|
 \le Ca\sum_i\gamma_i^2\delta_i^2.
\end{equation}

Use \eqref{F65:eq:bessel-absolute} when
$\Jfour\ge S^{9/2}$.  In the complementary regime,
\eqref{F65:eq:Jidentity} makes the radial segment from the uniform fan
quasiuniform; cyclic invariance kills the first derivative there, and
\eqref{F65:eq:bessel-hessian} with Taylor's formula gives
$CS\Jfour$.  This proves \eqref{F65:eq:bessel-relative}.
\end{proof}

\begin{corollary}[Correct-source Bessel gap]\label{F65:cor:bessel-gap}
There is $S_D>0$ such that, for $S\le S_D$,
\begin{equation}\label{F65:eq:bessel-gap}
 \boxed{
 P_D^G(\alpha)-P_D^G(\bar\alpha\mathbf1)
 \ge\frac{\zeta(3)}{40\pi^4}\Jfour(\alpha).}
\end{equation}
\end{corollary}

\begin{proof}
Subtract \eqref{F65:eq:bessel-relative} from \eqref{F65:eq:trace-gap} and decrease
$S_D$ until the relative error is at most $\zeta(3)/(40\pi^4)$.
\end{proof}

\subsection{The exact secant source}

Let
\begin{equation}\label{F65:eq:secant}
 A_\alpha=\sum_i\tan\frac{\alpha_i}{2},
 \qquad h_\alpha=\left(\frac\pi{A_\alpha}\right)^{1/2},
 \qquad
 \rho_\alpha^{\rm tan}=h_\alpha\sec x-1
 \quad\hbox{on }K_i.
\end{equation}
Writing $c_\alpha=h_\alpha-1$, one has
\begin{equation}\label{F65:eq:secant-split}
 \rho_\alpha^{\rm tan}=c_\alpha+G_\alpha+E_\alpha,
 \qquad
 E_\alpha=h_\alpha(\sec x-1)-\frac{x^2}{2}.
\end{equation}

\begin{lemma}[Scaled-spline secant remainder]\label{F65:lem:secant-remainder}
For every sufficiently fine positive fan,
\begin{equation}\label{F65:eq:secant-remainder}
 \|E_\alpha\|_{H^{1/2}}
 \le CS^2\|G_\alpha\|_{H^{1/2}}.
\end{equation}
The estimate is independent of the smallest cell.
\end{lemma}

\begin{proof}
On $K_i$, normalized by $x=(a+b)(y-t)/2$, both $G_\alpha$ and
$E_\alpha$ belong to a fixed compact family of even analytic profiles.
Taylor's formula gives, through one normalized derivative,
\begin{equation}\label{F65:eq:normalized-remainder}
 |E_\alpha|+(a+b)|E_\alpha'|
 \le CS^2(a+b)^2.
\end{equation}
The values from the two incident cells agree at each vertex.

For completeness, the mesh-independent assembly is as follows.  Split the
Slobodeckij integral into same-cell, adjacent-cell, and separated-cell
pieces.  Affine normalization and finite-dimensional norm equivalence apply
to the first two pieces.  For a separated pair $K_i,K_j$, subtract the cell
means and use
\begin{equation}
 \iint_{K_i\times K_j}\frac{d\theta\,d\varphi}
 {\operatorname {dist}_{\Sone}(\theta,\varphi)^2}
 \le C\frac{|K_i||K_j|}{d_{ij}^2}.
\end{equation}
Summing in dyadic annuli bounds the mean-difference energy by the same-cell
and adjacent-cell charges.  Equation \eqref{F65:eq:normalized-remainder} makes
every charge of $E_\alpha$ at most $CS^2$ times the corresponding charge of
$G_\alpha$.  This proves \eqref{F65:eq:secant-remainder} without a mesh-ratio
assumption.
\end{proof}

Define the exact tangential-source disk value
\begin{equation}\label{F65:eq:PDtan}
 P_D^{\rm tan}(\alpha)
 =\inf_{v\in W_\alpha}
 \qdisk\bigl(\PiNG(\rho_\alpha^{\rm tan}-c_\alpha)+v\bigr)
\end{equation}
and
\begin{equation}
 R_{\rm sec}(\alpha)=P_D^{\rm tan}(\alpha)-P_D^G(\alpha).
\end{equation}
Coercivity of \eqref{F65:eq:bessel-bounds}, the Hilbert distance inequality,
Lemma~\ref{F65:lem:secant-remainder}, and \eqref{F65:eq:Gbound} give
\begin{equation}\label{F65:eq:secant-absolute}
 |R_{\rm sec}(\alpha)|\le CS^5.
\end{equation}

\begin{lemma}[Relative secant correction]\label{F65:lem:secant-relative}
For every sufficiently fine positive fan,
\begin{equation}\label{F65:eq:secant-relative}
 |R_{\rm sec}(\alpha)-R_{\rm sec}(\bar\alpha\mathbf1)|
 \le C(S^{1/2}+S)\Jfour(\alpha).
\end{equation}
\end{lemma}

\begin{proof}
Work in a real quasiuniform material segment
$\gamma(t)=\gamma+t\delta$ and put
$h=\|\delta\|_{\ell^\infty}/a$.  Affine normalization of
\eqref{F65:eq:secant} and two differentiations of its elementary trigonometric
coefficients give, for $r=0,1,2$,
\begin{equation}
 \|\partial_t^rG_{\gamma(t)}\|_{H^{1/2}}
 \le Ch^ra^{3/2},
 \qquad
 \|\partial_t^rE_{\gamma(t)}\|_{H^{1/2}}
 \le Ch^ra^{7/2}.
\end{equation}
The real-material $q_D$ form and the exact trace/area graph have first and
second derivatives bounded by $Ch^r$ in their natural energy metrics; for
the Bessel part this is the order-minus-one kernel estimate used in
Proposition~\ref{F65:prop:bessel-relative}.  Differentiating the two coercive
Schur equations defining $P_D^{\rm tan}$ and $P_D^G$ therefore shows that
every second-derivative term in their difference is bounded by
\begin{equation}
 Ch^2(a^{3/2})(a^{7/2})=Ch^2a^5.
\end{equation}
Thus
\begin{equation}
 |D^2R_{\rm sec}(\gamma)[\delta,\delta]|
 \le Ca\sum_i\gamma_i^2\delta_i^2.
\end{equation}
The two-regime argument used in Proposition~\ref{F65:prop:bessel-relative},
with threshold $\Jfour=S^{9/2}$, proves \eqref{F65:eq:secant-relative}.
\end{proof}

\begin{theorem}[Exact tangential-polygon disk endpoint]\label{F65:thm:exact-disk}
There is $S_*>0$ such that, whenever $S\le S_*$,
\begin{equation}\label{F65:eq:exact-disk-gap}
 \boxed{
 P_D^{\rm tan}(\alpha)-P_D^{\rm tan}(\bar\alpha\mathbf1)
 \ge\frac{\zeta(3)}{80\pi^4}\Jfour(\alpha).}
\end{equation}
In particular, equality
$P_D^{\rm tan}(\alpha)=P_D^{\rm tan}(\bar\alpha\mathbf1)$ forces
$\alpha_i=\bar\alpha$ for every $i$.
\end{theorem}

\begin{proof}
Combine \eqref{F65:eq:bessel-gap} and \eqref{F65:eq:secant-relative}, decreasing
$S_*$ until the latter coefficient is at most
$\zeta(3)/(80\pi^4)$.  Strict convexity of $x^4$ gives
$\Jfour(\alpha)=0$ only at the uniform fan.
\end{proof}
\endgroup
